\documentclass[aspectratio=169,11pt]{beamer}
\usepackage[T2A]{fontenc}
\usepackage[utf8]{inputenc}
\usepackage[russian]{babel}
\usepackage{listings}
\usepackage{booktabs}
\usepackage{tabularx}
\usepackage{array}
\usepackage{tikz}
\usepackage{pgfplots}
\pgfplotsset{compat=1.16}
\usetikzlibrary{shapes.geometric,arrows.meta,positioning}

% ---------------------------------------------------------------------
% Тема и общие настройки
% ---------------------------------------------------------------------
\usetheme{Madrid}
\usecolortheme{default}
\setbeamertemplate{navigation symbols}{}
\setbeamertemplate{bibliography item}[text]
\setbeamerfont{normal text}{size=\normalsize}
\setbeamerfont{frametitle}{size=\Large}
\setbeamerfont{framesubtitle}{size=\large}
\setlength{\tabcolsep}{6pt}
\renewcommand{\arraystretch}{1.14}
\newcolumntype{Y}{>{\raggedright\arraybackslash}X}
\newcolumntype{C}{>{\centering\arraybackslash}X}
\newcolumntype{R}{>{\raggedleft\arraybackslash}X}

% ---------------------------------------------------------------------
% Макросы
% ---------------------------------------------------------------------
% Символ/идентификатор «как есть» (подчёркивания и т.п. писать без экранирования)
\newcommand{\sym}[1]{\texttt{\detokenize{#1}}}
% Прямые кавычки в моноширинном тексте (без активных символов babel)
\newcommand{\externC}{\texttt{extern~\textquotedbl C\textquotedbl}}
\newcommand{\todo}[1]{\textcolor{red}{[TODO: #1]}}
% Пустая ячейка в таблицах результатов
\newcommand{\tbd}{\textcolor{red}{\textbf{?}}}

% ---------------------------------------------------------------------
% Листинги
% ---------------------------------------------------------------------
\definecolor{deepblue}{rgb}{0,0,0.5}
\definecolor{deepred}{rgb}{0.6,0,0}
\definecolor{deepgreen}{rgb}{0,0.5,0}
\definecolor{halfgray}{rgb}{0.5,0.5,0.5}

\lstdefinelanguage{Rust}{
  morekeywords={as,async,await,break,const,continue,crate,dyn,else,enum,extern,false,fn,for,if,impl,in,let,loop,match,mod,move,mut,pub,ref,return,self,Self,static,struct,super,trait,true,type,unsafe,use,where,while,macro_rules,println},
  morecomment=[l]{//},
  morecomment=[s]{/*}{*/},
  morestring=[b]",
  sensitive=true
}

\lstdefinelanguage{Go}{
  morekeywords={break,case,chan,const,continue,default,defer,else,fallthrough,for,func,go,goto,if,import,interface,map,package,range,return,select,struct,switch,type,var},
  morecomment=[l]{//},
  morecomment=[s]{/*}{*/},
  morestring=[b]",
  sensitive=true
}

\lstdefinelanguage{Zig}{
  morekeywords={align,break,callconv,catch,comptime,const,continue,defer,else,enum,error,export,extern,fn,for,if,inline,pub,return,struct,switch,test,try,union,unreachable,var,while},
  morecomment=[l]{//},
  morestring=[b]",
  sensitive=true
}

\lstset{
  basicstyle=\ttfamily\small,
  keywordstyle=\bfseries\color{deepblue},
  emphstyle=\ttfamily\color{deepred},
  stringstyle=\color{deepgreen},
  commentstyle=\itshape\color{halfgray},
  breaklines=true,
  columns=fullflexible,
  keepspaces=true,
  showstringspaces=false,
  tabsize=2,
  numbers=left,
  numberstyle=\small\color{halfgray},
  numbersep=8pt,
  xleftmargin=2.4em,
  xrightmargin=0.6em,
  aboveskip=4pt,
  belowskip=4pt,
  rulesepcolor=\color{red!20!green!20!blue!20},
  frame=shadowbox,
  language=C++
}

% Примеры кода на слайдах: 13--15 строк, поэтому шрифт чуть мельче
\lstdefinestyle{code}{
  basicstyle=\ttfamily\footnotesize,
  numberstyle=\footnotesize\color{halfgray}
}

% Короткие команды оболочки: без номеров строк
\lstdefinestyle{shell}{
  language=bash,
  numbers=none,
  xleftmargin=0.6em,
  xrightmargin=0.6em
}

% ---------------------------------------------------------------------
% Титульные данные
% ---------------------------------------------------------------------
\title[Представление объектного файла]{Представление объектного файла: C++, Rust, Zig}
\subtitle{Форматы, символы, релокации, ABI и отладочная информация}
\author{Аржаев Дмитрий}
%\institute{\todo{Курс / кафедра}}
\date{29.09.2026}

\begin{document}

\begin{frame}
  \titlepage
\end{frame}

% =====================================================================
\section{Введение}

\begin{frame}{Зачем смотреть на объектный файл}
  \begin{columns}[T]
    \column{0.55\textwidth}
    \begin{itemize}
      \item Объектный файл («\texttt{.o}») --- результат компиляции одной единицы трансляции; вход линкера \cite[гл.~1--3]{levine}.
      \item Именно в нём фиксируется всё, что язык обещает внешнему миру: имена, ABI, режимы связывания.
      \item Три языка~--- C++, Rust, Zig: разные модели связывания и разная степень стандартизации ABI.
    \end{itemize}
    \column{0.45\textwidth}
    \centering
    \begin{tikzpicture}[node distance=4.5mm, every node/.style={font=\scriptsize},
        box/.style={draw, rounded corners, minimum width=30mm, minimum height=7mm, align=center},
        >=Stealth]
      \node[box] (src) {исходник};
      \node[box, below=of src] (cc) {компилятор};
      \node[box, below=of cc, fill=yellow!30] (obj) {объектный файл};
      \node[box, below=of obj] (ld) {линкер};
      \node[box, below=of ld] (exe) {исполняемый файл};
      \draw[->] (src) -- (cc);
      \draw[->] (cc) -- (obj);
      \draw[->] (obj) -- (ld);
      \draw[->] (ld) -- (exe);
    \end{tikzpicture}
  \end{columns}
\end{frame}

\begin{frame}{План и вопросы исследования}
  \begin{columns}[T]
    \column{0.48\textwidth}
    \begin{block}{План}
      \begin{enumerate}
        \item Обзор: форматы, секции, эволюция стандартов.
        \item Особенности C++/Rust/Zig на собственных примерах.
        \item Методология: метрики и стенд.
        \item Результаты сравнения.
      \end{enumerate}
    \end{block}
    \column{0.48\textwidth}
    \begin{block}{Вопросы}
      \begin{itemize}
        \item Как различаются размер и состав \texttt{.o}?
        \item Сколько работы остаётся линкеру?
        \item Насколько совместим каждый язык с C~ABI?
        \item Воспроизводим ли результат?
      \end{itemize}
    \end{block}
  \end{columns}
\end{frame}

% =====================================================================
\section{Обзор}
% [~8 min]

\begin{frame}{Форматы объектных файлов: эволюция и сравнение}
  \begin{itemize}
    \item Линия развития: \texttt{a.out} $\rightarrow$ COFF $\rightarrow$ ELF; PE выросла из COFF; Mach-O~--- формат семейства NeXTSTEP/macOS/iOS \cite{levine,pecoff,macho}.
    \item Объектный файл на разных ОС: ELF (\texttt{ET\_REL}) \cite{tis}, COFF (\texttt{.obj}) \cite{pecoff}, Mach-O (\texttt{MH\_OBJECT}) \cite{macho}.
  \end{itemize}
  \small
  \begin{center}
    \begin{tabularx}{0.90\textwidth}{lCCC}
      \toprule
                & ELF \cite{tis}       & PE/COFF \cite{pecoff} & Mach-O \cite{macho}        \\
      \midrule
      ОС        & Linux, BSD, Solaris  & Windows               & macOS, iOS                 \\
      Единица   & секции               & секции                & сегменты $\to$ секции      \\
      Имя кода  & \texttt{.text}       & \texttt{.text}        & \texttt{\_\_TEXT,\_\_text} \\
      Каталог   & таблица секций       & таблица секций        & load commands              \\
      \bottomrule
    \end{tabularx}
  \end{center}
\end{frame}

\begin{frame}{Структура relocatable ELF (\texttt{.o})}
  \begin{columns}[T]
    \column{0.42\textwidth}
    \centering
    \begin{tikzpicture}[every node/.style={font=\scriptsize, draw, minimum width=42mm, minimum height=6.5mm}]
      \node[fill=blue!10] (h) {ELF header};
      \node[below=0mm of h] (t) {\texttt{.text}};
      \node[below=0mm of t] (d) {\texttt{.data} / \texttt{.bss} / \texttt{.rodata}};
      \node[below=0mm of d] (s) {\texttt{.symtab} / \texttt{.strtab}};
      \node[below=0mm of s] (r) {\texttt{.rela.*}};
      \node[below=0mm of r] (g) {\texttt{.debug\_*} / \texttt{.eh\_frame}};
      \node[below=0mm of g, fill=blue!10] (sh) {таблица заголовков секций};
    \end{tikzpicture}
    \column{0.58\textwidth}
    \begin{itemize}
      \item ELF header: класс (32/64), порядок байт, машина, тип файла \cite{tis}.
      \item Таблица секций описывает содержимое; для \texttt{.o} она главная \cite{tis}.
      \item Таблица сегментов (program headers) нужна загрузчику; у relocatable-файла обычно отсутствует \cite{tis}.
      \item \texttt{.bss} (NOBITS) не занимает места в файле.
      \item Подробнее: \texttt{man 5 elf}.
    \end{itemize}
  \end{columns}
\end{frame}

\begin{frame}{Основные секции}
  \small
  \begin{center}
    \begin{tabularx}{0.90\textwidth}{lY}
      \toprule
      Секция \cite{tis}                                   & Назначение                                                  \\
      \midrule
      \texttt{.text}                                      & машинный код                                                \\
      \texttt{.data}, \texttt{.bss}                       & инициализированные / нулевые изменяемые данные              \\
      \texttt{.rodata}                                    & константы, строковые литералы                               \\
      \texttt{.symtab}, \texttt{.strtab}                  & таблица символов и строк                                    \\
      \texttt{.rela.text}, \texttt{.rela.*}               & записи релокаций                                            \\
      \texttt{.eh\_frame}, \texttt{.gcc\_except\_table}   & раскрутка стека и обработка исключений \cite{itanium}       \\
      \texttt{.group} / COMDAT                            & дедупликация inline-функций и шаблонов \cite{itanium}       \\
      \texttt{.debug\_*}                                  & отладочная информация DWARF \cite{dwarf}                    \\
      \texttt{.note.GNU-stack}                            & признак исполняемости стека \cite{maskray}                  \\
      \bottomrule
    \end{tabularx}
  \end{center}
\end{frame}

\begin{frame}{Что стандарты говорят про объектный файл}
  \begin{itemize}
    \item \textbf{C++}: стандарт ISO/IEC 14882 описывает фазы трансляции, linkage и ODR, но не формат \texttt{.o}. Формат и ABI задают Itanium C++ ABI, System~V psABI и MSVC \cite{iso,itanium,psabi}.
    \item \textbf{Rust}: ISO-стандарта нет; поведение описано в Rust Reference и RFC (например, RFC~2603 про mangling) \cite{rustref,rfc2603}.
    \item \textbf{Zig}: язык описан в Language Reference; до версии 1.0 стабильность ABI и формата не гарантируется \cite{zigdoc}.
  \end{itemize}
  \bigskip
  \alert{Вывод:} единственный «общий язык» между тремя компиляторами~--- C~ABI платформы.
\end{frame}

\begin{frame}{Как развитие C++ влияло на объектный файл}
  \small
  \begin{center}
    \begin{tabularx}{0.96\textwidth}{lYY}
      \toprule
      Стандарт & Что появилось                                & Следствие для \texttt{.o}                                                        \\
      \midrule
      C++98    & ODR, \externC, шаблоны, inline \cite{cppref} & mangling, weak/COMDAT-символы \cite{itanium}                                     \\
      C++11    & \texttt{thread\_local}                       & TLS-секции и TLS-релокации \cite{drepperTLS}                                     \\
      C++17    & inline-переменные \cite{cppref}              & COMDAT-группы для данных \cite{itanium}                                          \\
      C++20    & модули \cite{cppref}                         & BMI как новый артефакт рядом с \texttt{.o}; формат не стандартизован \cite{cppref} \\
      \bottomrule
    \end{tabularx}
  \end{center}
\end{frame}

% =====================================================================
\section{Особенности языков}
% [~8 min]

\begin{frame}{Символы в ELF}
  \begin{itemize}
    \item Запись символа: имя, значение, размер, тип, binding, visibility, индекс секции \cite{tis}.
    \item Binding: \texttt{LOCAL}, \texttt{GLOBAL}, \texttt{WEAK}. Тип: \texttt{FUNC}, \texttt{OBJECT}, \texttt{TLS}, \texttt{SECTION}, \texttt{FILE} \cite{tis}.
    \item Индекс \texttt{UND}~--- символ не определён в этом файле; его разрешит линкер \cite{tis}.
  \end{itemize}
  \medskip
  Просмотр: \texttt{readelf -sW file.o}, \texttt{nm -C file.o}, \texttt{llvm-readobj -{}-symbols file.o}.
\end{frame}

\begin{frame}[fragile]{Свой пример: C++}
\begin{lstlisting}[style=code, language=C++]
namespace geom {
  int area(int w, int h) {
    return w * h;
  }
}
template <class T>
T add(T a, T b) {
  return a + b;
}
template int add<int>(int, int);
extern "C" int c_area(int w, int h) {
  return geom::area(w, h);
}
\end{lstlisting}
  \smallskip
  \centering
  \textit{Смотрим символы: \texttt{geom::area}, \texttt{add<int>}, \texttt{c\_area}.}
\end{frame}

\begin{frame}[fragile]{Свой пример: Rust 2024}
\begin{lstlisting}[style=code, language=Rust, emph={i32,T}]
pub mod geom {
  pub fn area(w: i32, h: i32) -> i32 {
    w * h
  }
}
pub fn add<T: core::ops::Add<Output = T>>(a: T, b: T) -> T {
  a + b
}
pub fn add_i32(a: i32, b: i32) -> i32 {
  add(a, b)
}
#[unsafe(no_mangle)]
pub extern "C" fn c_area(w: i32, h: i32) -> i32 {
  geom::area(w, h)
}
\end{lstlisting}
  \smallskip
  \centering
  \textit{Смотрим символы: \texttt{geom::area}, \texttt{add\_i32} (в него подставлен \texttt{add::<i32>}), \texttt{c\_area}.}
\end{frame}

\begin{frame}[fragile]{Свой пример: Zig}
\begin{lstlisting}[style=code, language=Zig, emph={i32,type,T}]
const geom = struct {
  fn area(w: i32, h: i32) i32 {
    return w * h;
  }
};
fn add(comptime T: type, a: T, b: T) T {
  return a + b;
}
export fn add_i32(a: i32, b: i32) i32 {
  return add(i32, a, b);
}
export fn c_area(w: i32, h: i32) i32 {
  return geom.area(w, h);
}
\end{lstlisting}
  \smallskip
  \centering
  \textit{Экспортируются только \texttt{add\_i32} и \texttt{c\_area}; \texttt{geom.area} и \texttt{add} остаются внутренними.}
\end{frame}

\begin{frame}[fragile]{Как смотреть символы и релокации}
  \begin{columns}[T,onlytextwidth]
    \column{0.48\textwidth}
    \textbf{Символы}
    \begin{lstlisting}[style=shell]
readelf -sW file.o
nm -C file.o
llvm-readobj --symbols file.o
    \end{lstlisting}
    \column{0.48\textwidth}
    \textbf{Релокации}
    \begin{lstlisting}[style=shell]
readelf -rW file.o
objdump -dr file.o
    \end{lstlisting}
  \end{columns}
  \bigskip
  Инструменты: GNU binutils \cite{binutils} и LLVM \cite{llvmtools}.
\end{frame}

\begin{frame}{Mangling: C++ (Itanium ABI)}
  \begin{itemize}
    \item Имя кодирует область видимости, имя и типы параметров: префикс \sym{_Z}, вложенные имена в \sym{N...E} \cite{itanium}.
    \item \externC{} отключает mangling: остаётся имя \sym{c_area} \cite{itanium}.
  \end{itemize}
  \small
  \begin{center}
    \begin{tabularx}{0.80\textwidth}{lY}
      \toprule
      Сущность                                & Символ                             \\
      \midrule
      \texttt{geom::area(int, int)}           & \sym{_ZN4geom4areaEii}             \\
      \texttt{add<int>(int, int)}             & \sym{_Z3addIiET_S0_S0_}            \\
      \texttt{c\_area} (\externC)             & \sym{c_area}                       \\
      vtable / typeinfo класса                & префиксы \sym{_ZTV}, \sym{_ZTI}    \\
      \bottomrule
    \end{tabularx}
  \end{center}
  \medskip
  Проверка: \texttt{c++filt}, \texttt{nm} на \texttt{.o}, Compiler Explorer \cite{godbolt}.
\end{frame}

\begin{frame}{Mangling: Rust и Zig}
  \begin{itemize}
    \item \textbf{Rust v0} (RFC 2603): префикс \sym{_R} и структурированное кодирование пути и параметров символа \cite{rfc2603}. Стабильный \texttt{rustc} использует эту схему по умолчанию начиная с Rust~1.97; версию \texttt{rustc} нужно зафиксировать в эксперименте \cite{rustcbook}.
    \item \textbf{Rust legacy}: историческая схема с именами вида \sym{_ZN4demo4geom4area17h<16 hex>E}; до Rust~1.97 была схемой по умолчанию \cite{rustcbook}.
    \item \textbf{Rust FFI}: \texttt{\#[unsafe(no\_mangle)]} и \externC{} --- имя экспортируемого символа задаётся явно \cite{rustref}.
    \item \textbf{Zig}: \texttt{export fn} даёт символ с исходным именем (\sym{c_area}); неэкспортируемые функции остаются внутренними \cite{zigdoc}.
  \end{itemize}
\end{frame}

\begin{frame}{Релокации}
  \small
  \begin{itemize}
    \item Компилятор не знает финальных адресов, поэтому оставляет в \texttt{.rela.*} записи: смещение, тип, символ, addend \cite{tis}.
    \item Линкер вычисляет значение по формуле типа \cite{psabi}: $S$~--- адрес символа, $A$~--- addend, $P$~--- изменяемое место, $L$~--- запись в PLT, $G$~--- запись в GOT, $GOT$~--- адрес GOT.
  \end{itemize}
  \footnotesize
  \begin{center}
    \renewcommand{\arraystretch}{1.0}
    \setlength{\tabcolsep}{4pt}
    \begin{tabularx}{0.98\textwidth}{lcY}
      \toprule
      Тип (x86-64)                  & Формула      & Где встречается                                      \\
      \midrule
      \texttt{R\_X86\_64\_64}       & $S+A$        & 64-битный абсолютный адрес \cite{psabi}              \\
      \texttt{R\_X86\_64\_PC32}     & $S+A-P$      & доступ к данным относительно PC \cite{bender}        \\
      \texttt{R\_X86\_64\_PLT32}    & $L+A-P$      & вызов внешней функции \cite{bender}                  \\
      \texttt{R\_X86\_64\_GOTPCREL} & $G+GOT+A-P$  & доступ через GOT \cite{bender}                       \\
      \texttt{R\_X86\_64\_TPOFF32}  & TLS-смещение & thread-local в исполняемом файле \cite{drepperTLS}   \\
      \bottomrule
    \end{tabularx}
  \end{center}
  Все три языка собираются через LLVM-бэкенд, поэтому наборы релокаций близки, но не обязаны совпадать. Их число зависит от PIC/PIE, TLS и отладочной информации (в \texttt{.rela.debug\_*} их много) \cite{maskray}.
\end{frame}

\begin{frame}{ABI и calling conventions}
  \small
  \begin{center}
    \begin{tabularx}{0.86\textwidth}{lCC}
      \toprule
                                & System V AMD64 \cite{psabi}   & Microsoft x64 \cite{pecoff} \\
      \midrule
      Целочисленные аргументы   & rdi, rsi, rdx, rcx, r8, r9    & rcx, rdx, r8, r9            \\
      Вещественные аргументы    & xmm0--xmm7                    & xmm0--xmm3                  \\
      Возврат                   & rax (rdx), xmm0               & rax, xmm0                   \\
      Shadow space              & нет                           & 32 байта                    \\
      \bottomrule
    \end{tabularx}
  \end{center}
  \begin{itemize}
    \item \textbf{C++}: соглашение платформы; на Linux с GCC/Clang layout классов и vtable задаёт Itanium C++ ABI, на Windows используется ABI MSVC \cite{itanium}.
    \item \textbf{Rust}: \texttt{extern~\textquotedbl Rust\textquotedbl} --- ABI не является стабильным контрактом. Для FFI используют \externC{} (ABI функций) и \texttt{\#[repr(C)]} (layout типов) \cite{rustref}.
    \item \textbf{Zig}: соглашение задаётся через \texttt{callconv} (\texttt{.C} в старых версиях, \texttt{.c} в новых); \texttt{export} использует C~ABI \cite{zigdoc}.
  \end{itemize}
\end{frame}

\begin{frame}{Исключения, RTTI и метаданные}
  \small
  \begin{center}
    \begin{tabularx}{0.98\textwidth}{lYYY}
      \toprule
                         & C++ \cite{itanium}            & Rust \cite{rustref}                                & Zig \cite{zigdoc}       \\
      \midrule
      Ошибки             & исключения                    & \texttt{Result}; \texttt{panic} (unwind или abort) & error union (значение)  \\
      Раскрутка стека    & \texttt{.eh\_frame}, LSDA     & \texttt{.eh\_frame} при \texttt{panic=unwind}      & не требуется для ошибок \\
      RTTI               & typeinfo, vtable-символы      & нет аналога; \texttt{TypeId} на уровне языка       & нет                     \\
      Метаданные модулей & BMI (C++20), вне \texttt{.o}  & метаданные крейта (\texttt{rlib}), вне \texttt{.o} & нет                     \\
      \bottomrule
    \end{tabularx}
  \end{center}
  \medskip
  Практическая проверка: сравнить \texttt{.eh\_frame}, \texttt{.gcc\_except\_table} и число символов \sym{_ZTV}/\sym{_ZTI} на одном коде с \texttt{-fno-exceptions} / \texttt{panic=abort}.
\end{frame}

\begin{frame}{Отладочная информация}
  \begin{itemize}
    \item ELF: DWARF (v5 в современных компиляторах): \texttt{.debug\_info}, \texttt{.debug\_abbrev}, \texttt{.debug\_line}, \texttt{.debug\_str}, \texttt{.debug\_line\_str}, \texttt{.debug\_rnglists} \cite{dwarf}.
    \item Windows: CodeView/PDB. macOS: DWARF в объектных файлах, сборка в dSYM.
    \item Управление: \texttt{-g}, \texttt{-gsplit-dwarf}, сжатие секций; Rust: \texttt{-C debuginfo=0|1|2} \cite{rustcbook}; Zig: режим \texttt{Debug} содержит отладочную информацию \cite{zigdoc}.
    \item В \texttt{.o} DWARF сильно нагружает линкер: релокации ссылаются на \texttt{.text} и строки.
  \end{itemize}
\end{frame}

% =====================================================================
\section{Методология}
% [~4 min]

\begin{frame}{Метрики: объектный файл}
  \small
  \begin{center}
    \begin{tabularx}{0.98\textwidth}{lYY}
      \toprule
      Метрика                    & Инструмент                                          & Пояснение                                                                                  \\
      \midrule
      Размер \texttt{.o} и по секциям & \texttt{size -A}, \texttt{readelf -SW}         & total, \texttt{.text}, \texttt{.data}, \texttt{.rodata}, \texttt{.debug\_*}, \texttt{.eh\_frame} \\
      Число секций               & \texttt{readelf -SW}                                & ---                                                                                        \\
      Символы                    & \texttt{nm}, \texttt{readelf -sW}                   & defined / undefined; local / global / weak                                                 \\
      Релокации                  & \texttt{readelf -rW}                                & число и типы; с \texttt{-g} и без                                                          \\
      \bottomrule
    \end{tabularx}
  \end{center}
  \smallskip
  Сложность ABI (mangling, исключения, RTTI, метаданные) не сводится к одному числу и представлена качественной таблицей. Инструменты: binutils \cite{binutils}, LLVM \cite{llvmtools}, заметки \cite{maskray}.
\end{frame}

\begin{frame}{Метрики: сборка и линковка}
  \small
  \begin{center}
    \begin{tabularx}{0.98\textwidth}{lYY}
      \toprule
      Метрика                       & Инструмент                                              & Пояснение                    \\
      \midrule
      Время и память компиляции     & \texttt{hyperfine} \cite{hyperfine}, \texttt{/usr/bin/time -v} & медиана и разброс, Max RSS \\
      Работа линкера                & размер после \texttt{-{}-gc-sections}, LTO \cite{taylor} & undefined, релокации, DCE    \\
      Совместимость с C~ABI         & чек-лист                                                & вызов из C, вызов C, стабильность \\
      Воспроизводимость             & \texttt{sha256sum}                                      & две сборки в разных каталогах \\
      \bottomrule
    \end{tabularx}
  \end{center}
  \smallskip
  Метрики фиксируют как свойства объектного файла, так и стоимость сборки и последующей линковки.
\end{frame}

\begin{frame}{Стенд и корпус}
  \begin{itemize}
    \item Платформа: x86-64 Linux, ELF; \texttt{clang++ 22.1.8}, \texttt{rustc 1.98.1}, \texttt{zig 0.16.0}; \texttt{readelf 2.47}.
    \item В сырых данных присутствуют режимы \texttt{O0g}, \texttt{O2}, \texttt{O2g}, \texttt{Os}; таблицы \texttt{report.py} по умолчанию строятся только для \texttt{O2} и \texttt{O2g}, а графики --- для \texttt{O2g}.
    \begin{itemize}
      \item C++: \texttt{-O0}, \texttt{-O2}, \texttt{-Os};
      \item Rust: \texttt{-C opt-level=0|2|s} \cite{rustcbook};
      \item Zig: \texttt{Debug}, \texttt{ReleaseFast}, \texttt{ReleaseSmall} \cite{zigdoc}.
    \end{itemize}
    \item Корпус: 5 программ на каждый язык: \texttt{demo}, \texttt{generics}, \texttt{strings}, \texttt{errors}, \texttt{cabi}.
  \end{itemize}
\end{frame}

\begin{frame}[fragile]{Команды сборки}
\begin{lstlisting}[style=shell]
# C++
clang++ -c -O2 -g src.cpp -o out.o

# Rust
rustc --crate-type=lib --emit=obj -C opt-level=2 \
  -C debuginfo=2 -C codegen-units=1 src.rs -o out.o

# Zig
zig build-obj src.zig -O ReleaseFast -fllvm \
  -femit-bin=out.o
\end{lstlisting}
  \smallskip
  Для честности сравнения в стенде используется LLVM-бэкенд у всех трёх (в Zig явно \texttt{-fllvm}) \cite{zigdoc}; у Rust задано \texttt{codegen-units=1}, чтобы получить один \texttt{.o} \cite{rustcbook}. Это ограничение эксперимента, а не свойство языков.
\end{frame}

\begin{frame}{Воспроизводимость}
  \begin{itemize}
    \item Собираем один и тот же исходник дважды: в разных каталогах и в разное время \cite{reprod}.
    \item Сравниваем \texttt{sha256} объектных файлов; при отличии смотрим \texttt{diff} вывода \texttt{readelf}/\texttt{objdump}.
    \item Типичные источники различий: абсолютные пути в DWARF, метки времени, хэши в символах \cite{reprod}.
    \item Инструменты: \texttt{-ffile-prefix-map} (C++) \cite{reprod}, \texttt{-{}-remap-path-prefix} (Rust) \cite{rustcbook}.
  \end{itemize}
\end{frame}

% =====================================================================
\section{Результаты}

\begin{frame}{Размер объектного файла}
  \small
  \begin{center}
    \setlength{\tabcolsep}{4pt}
    \IfFileExists{../results/table_sizes.tex}{%
      % generated by scripts/report.py, do not edit manually
\begin{tabular}{@{}lrrrrr@{}}
\toprule
Language / mode & total & \texttt{.text} & \texttt{.data+.rodata} & \texttt{.debug\_*} & \texttt{.eh\_frame} \\
\midrule
C++ O2 & 14\,232 & 3\,653 & 164 & 0 & 720 \\
C++ O2g & 53\,640 & 3\,653 & 164 & 21\,344 & 720 \\
Rust O2 & 14\,592 & 2\,250 & 349 & 0 & 560 \\
Rust O2g & 82\,104 & 2\,250 & 349 & 33\,998 & 560 \\
Zig O2 & 311\,256 & 105\,851 & 43\,793 & 0 & 9\,432 \\
Zig O2g & 4\,311\,752 & 333\,163 & 63\,637 & 2\,182\,518 & 13\,600 \\
\bottomrule
\end{tabular}
%
    }{%
      \fbox{\parbox{0.9\textwidth}{\centering
        Не найден \texttt{results/table\_sizes.tex}. Сначала запустите
        \texttt{scripts/report.py} (по умолчанию: режимы \texttt{O2,O2g}).}}}%
  \end{center}

  \smallskip
  \scriptsize
  Все числовые значения в таблице суммированы по корпусу из 5 программ на каждый язык.
  График \texttt{chart\_sizes.png} скрипт строит отдельно для режимов
  \texttt{O2} и \texttt{O2g}.
\end{frame}

\begin{frame}{Символы, секции и релокации}
  \small
  \begin{center}
    \setlength{\tabcolsep}{4pt}
    \IfFileExists{../results/table_symbols.tex}{%
      % generated by scripts/report.py, do not edit manually
\begin{tabular}{@{}lrrrrl@{}}
\toprule
Language / mode & sections & def./undef. & weak & relocations & main type \\
\midrule
C++ O2 & 69 & 55/11 & 8 & 56 & \texttt{R\_X86\_64\_PC32} \\
C++ O2g & 132 & 95/11 & 8 & 571 & \texttt{R\_X86\_64\_32} \\
Rust O2 & 83 & 62/7 & 2 & 53 & \texttt{R\_X86\_64\_PC32} \\
Rust O2g & 141 & 91/7 & 2 & 1\,242 & \texttt{R\_X86\_64\_32} \\
Zig O2 & 53 & 168/7 & 0 & 5\,878 & \texttt{R\_X86\_64\_PC32} \\
Zig O2g & 112 & 765/8 & 1 & 69\,823 & \texttt{R\_X86\_64\_64} \\
\bottomrule
\end{tabular}
%
    }{%
      \fbox{\parbox{0.9\textwidth}{\centering
        Не найден \texttt{results/table\_symbols.tex}. Сначала запустите
        \texttt{scripts/report.py}.}}}%
  \end{center}

  \smallskip
  \scriptsize
  Учитывается общее число релокаций; в отчёте дополнительно выводится
  число отладочных релокаций. Основной тип --- наиболее частый тип
  релокации для данной пары язык/режим.
\end{frame}

\begin{frame}{Время и память компиляции; работа линкера}
  \small
  \begin{center}
    \setlength{\tabcolsep}{4pt}
    \IfFileExists{../results/table_time.tex}{%
      % generated by scripts/report.py, do not edit manually
\begin{tabular}{@{}lrrrrr@{}}
\toprule
Language / mode & time, ms & Max RSS, MB & \texttt{.text*} sections & COMDAT groups & undef. \\
\midrule
C++ O2 & 535 & 177 & 12 & 8 & 11 \\
C++ O2g & 546 & 178 & 12 & 8 & 11 \\
Rust O2 & 224 & 188 & 21 & 2 & 7 \\
Rust O2g & 247 & 189 & 21 & 2 & 7 \\
Zig O2 & 4\,457 & 214 & 6 & 0 & 7 \\
Zig O2g & 13\,615 & 322 & 6 & 0 & 8 \\
\bottomrule
\end{tabular}
%
    }{%
      \fbox{\parbox{0.9\textwidth}{\centering
        Не найден \texttt{results/table\_time.tex}. Сначала запустите
        \texttt{scripts/report.py}.}}}%
  \end{center}

  \smallskip
  \scriptsize
  Время компиляции --- сумма медиан по программам корпуса;
  Max RSS --- максимум по программам. `.text*` sections и COMDAT groups
  показывают число единиц, которые могут быть затронуты
  \texttt{-{}-gc-sections}.
\end{frame}

\begin{frame}{Сравнение по ABI и совместимости (качественно)}
  \footnotesize
  \begin{center}
    \begin{tabularx}{0.98\textwidth}{lYYY}
      \toprule
                                 & C++ \cite{itanium}                      & Rust \cite{rustref}                       & Zig \cite{zigdoc}                    \\
      \midrule
      Mangling                   & Itanium ABI (Linux/Clang)               & v0 (по умолчанию с Rust~1.97)             & без mangling для \texttt{export}     \\
      Стабильный ABI языка       & нет (ABI задают платформа и реализация)   & нет; для FFI используют C~ABI             & нет                                  \\
      Вызов из C без обёрток     & \externC                                & \externC{} + \texttt{no\_mangle}          & \texttt{export}                      \\
      Исключения / раскрутка     & есть                                    & \texttt{panic=unwind}                     & нет для ошибок                       \\
      RTTI-символы               & есть                                    & нет                                       & нет                                  \\
      Метаданные вне \texttt{.o} & BMI (модули)                            & \texttt{rlib}-метаданные                  & нет                                  \\
      \bottomrule
    \end{tabularx}
  \end{center}
  \medskip
  Таблицу нужно проверить на выбранных версиях компиляторов.
\end{frame}

\begin{frame}{Выводы}
  \begin{itemize}

    \item \textbf{Размер \texttt{.o}.}
    В режиме \texttt{O2} C++ и Rust дают близкий суммарный размер:
    14\,232 и 14\,592 байт, тогда как Zig --- 311\,256 байт.
    При \texttt{O2g} размер возрастает до 53\,640, 82\,104 и
    4\,311\,752 байт соответственно; значительная часть этого роста
    приходится на \texttt{.debug\_\*}.

    \item \textbf{Секции, символы и релокации.}
    При добавлении отладочной информации число секций увеличивается
    у всех трёх языков: с 69 до 132 у C++, с 83 до 141 у Rust
    и с 53 до 112 у Zig.
    Число релокаций возрастает значительно сильнее:
    56$\rightarrow$571, 53$\rightarrow$1\,242 и
    5\,878$\rightarrow$69\,823.

    \item \textbf{Стоимость компиляции.}
    Переход от \texttt{O2} к \texttt{O2g} почти не меняет время
    и Max RSS у C++ и Rust
    (535$\rightarrow$546 мс и 177$\rightarrow$178 МБ;
    224$\rightarrow$247 мс и 188$\rightarrow$189 МБ).
    У Zig рост заметнее: 4\,457$\rightarrow$13\,615 мс
    и 214$\rightarrow$322 МБ.

  \end{itemize}
\end{frame}

\begin{frame}{Выводы}
  \begin{itemize}
    \item \textbf{Структура результатов.}
    У C++ и Rust число COMDAT-групп составляет 8 и 2 соответственно
    и не меняется между режимами; у Zig в этих данных COMDAT-группы
    отсутствуют.
    Для \texttt{O2g} основным типом релокации является
    \texttt{R\_X86\_64\_32} для C++ и Rust и
    \texttt{R\_X86\_64\_64} для Zig.
  \end{itemize}
    \begin{alertblock}{Ограничения}
    Результаты относятся к данному корпусу программ, конкретным версиям
    компиляторов и платформе x86-64 Linux; таблицы отражают суммарные
    значения по корпусу.
  \end{alertblock}
\end{frame}


\begin{frame}{Заключение}
  \begin{itemize}
    \item Объектный файл~--- точка, где встречаются язык, компилятор, ABI и линкер.
    \item Общий знаменатель трёх языков~--- C~ABI; всё остальное определяется реализацией.
    \item Исходники, скрипты сбора метрик, сырые данные доступны в \textbf{\href{https://github.com/InfernalKn1ght/cpp-object-files-talk}{репозитории}}.
  \end{itemize}
  \vfill
\end{frame}

% =====================================================================
\begin{frame}[allowframebreaks]{Источники}
  \small
  \begin{thebibliography}{99}
    \bibitem{tis} TIS Committee. \emph{Tool Interface Standard (TIS) Executable and Linking Format (ELF) Specification}, v1.2, 1995.
    \bibitem{levine} J.~R.~Levine. \emph{Linkers and Loaders}. Morgan Kaufmann, 2000.
    \bibitem{psabi} System V Application Binary Interface, AMD64 Architecture Processor Supplement (x86-64 psABI). \texttt{gitlab.com/x86-psABIs}.
    \bibitem{itanium} Itanium C++ ABI. \texttt{itanium-cxx-abi.github.io/cxx-abi}.
    \bibitem{pecoff} Microsoft. PE Format; x64 calling convention. \texttt{learn.microsoft.com}.
    \bibitem{macho} Apple. \emph{Mach-O File Format Reference} (архивная документация).
    \bibitem{dwarf} DWARF Debugging Information Format, Version 5. \texttt{dwarfstd.org}.
    \bibitem{iso} ISO/IEC 14882 --- Programming languages --- C++.
    \bibitem{cppref} cppreference.com: inline specifier, linkage, modules.
    \bibitem{taylor} I.~L.~Taylor. Linkers (серия статей). \texttt{airs.com/ian}, 2007--2008.
    \bibitem{drepperTLS} U.~Drepper. \emph{ELF Handling For Thread-Local Storage}, 2013.
    \bibitem{bender} E.~Bendersky. Position Independent Code (PIC) in shared libraries on x64. \texttt{eli.thegreenplace.net}.
    \bibitem{maskray} F.~Song. MaskRay blog. \texttt{maskray.me}.
    \bibitem{rfc2603} Rust RFC 2603: Symbol Name Mangling v0.
    \bibitem{rustref} The Rust Reference (type layout, ABI, linkage).
    \bibitem{rustcbook} The rustc book: codegen options, symbol mangling.
    \bibitem{zigdoc} Zig Language Reference. \texttt{ziglang.org/documentation}.
    \bibitem{reprod} Reproducible Builds: documentation. \texttt{reproducible-builds.org}.
    \bibitem{godbolt} Compiler Explorer. \texttt{godbolt.org}.
    \bibitem{binutils} GNU Binutils documentation. \texttt{sourceware.org/binutils/docs}.
    \bibitem{llvmtools} LLVM Command Guide (\texttt{llvm-readobj} и др.). \texttt{llvm.org/docs/CommandGuide}.
    \bibitem{hyperfine} D.~Peter. hyperfine: command-line benchmarking tool. \texttt{github.com/sharkdp/hyperfine}.
  \end{thebibliography}
\end{frame}

\end{document}